\section{Splitting fields and the structure of group algebras}
\begin{definition}
    Let \(A\) be a finite-dimensional \(k\)-algebra. We say that \(A\) is
    \emph{split} over \(k\) if for every simple left \(A\)-module \(S\), one has
    \[
        \operatorname{End}_A(S)=k\operatorname{id}_S.
    \]
    Equivalently, every simple \(A\)-module has endomorphism ring exactly \(k\).

    If \(A=kG\), then we say that \(k\) is a \emph{splitting field} for \(G\) if
    \(kG\) is split over \(k\).
\end{definition}

\begin{theorem}[Finite-dimensional density theorem]\label{thm:finite-dimensional-density}
    Let \(A\) be a \(k\)-algebra, and let \(S\) be a finite-dimensional simple
    left \(A\)-module. Put
    \[
        D:=\operatorname{End}_A(S)^{\operatorname{op}}.
    \]
    We regard \(S\) as a right \(D\)-vector space by
    \[
        s\cdot \varphi := \varphi(s),
        \qquad
        s\in S,\ \varphi\in \operatorname{End}_A(S).
    \]
    Let
    \[
        I=\operatorname{Ann}_A(S)
        =
        \{\,a\in A\mid as=0\text{ for all }s\in S\,\}.
    \]
    Then the natural algebra homomorphism
    \[
        \overline{\rho}:A/I\longrightarrow \operatorname{End}_D(S),
        \qquad
        a+I\longmapsto L_a,
    \]
    where
    \[
        L_a:S\longrightarrow S,\qquad s\longmapsto as,
    \]
    is an isomorphism.
\end{theorem}

\begin{proof}
    By Schur's lemma,
    \[
        \operatorname{End}_A(S)
    \]
    is a division ring. Hence
    \[
        D=\operatorname{End}_A(S)^{\operatorname{op}}
    \]
    is also a division ring, and \(S\) is a right vector space over \(D\).

    For every \(a\in A\), the map
    \[
        L_a:S\longrightarrow S,\qquad s\longmapsto as,
    \]
    is right \(D\)-linear. Indeed, for \(s\in S\) and
    \(\varphi\in \operatorname{End}_A(S)\), we have
    \[
        L_a(s\cdot \varphi)
        =
        L_a(\varphi(s))
        =
        a\varphi(s)
        =
        \varphi(as)
        =
        L_a(s)\cdot \varphi,
    \]
    because \(\varphi\) is \(A\)-linear. Thus we get an algebra homomorphism
    \[
        \rho:A\longrightarrow \operatorname{End}_D(S),
        \qquad
        a\longmapsto L_a.
    \]

    Its kernel is precisely
    \[
        \operatorname{Ker}(\rho)
        =
        \{\,a\in A\mid L_a=0\,\}
        =
        \{\,a\in A\mid as=0\text{ for all }s\in S\,\}
        =
        \operatorname{Ann}_A(S)
        =
        I.
    \]
    Therefore \(\rho\) induces an injective homomorphism
    \[
        \overline{\rho}:A/I\hookrightarrow \operatorname{End}_D(S).
    \]

    It remains to prove that \(\overline{\rho}\) is surjective.

    We first prove the following density claim. Let
    \[
        x_1,\dots,x_m\in S
    \]
    be right \(D\)-linearly independent, and let
    \[
        y_1,\dots,y_m\in S
    \]
    be arbitrary. Then there exists \(a\in A\) such that
    \[
        ax_i=y_i
        \qquad
        (1\leq i\leq m).
    \]

    To prove the claim, consider the \(A\)-submodule
    \[
        N
        :=
        \{\, (ax_1,\dots,ax_m)\in S^m\mid a\in A\,\}
        \subseteq S^m.
    \]
    We claim that
    \[
        N=S^m.
    \]
    Suppose not. Since \(S^m\) is a finite direct sum of copies of the simple
    module \(S\), it is semisimple. Hence the proper submodule \(N\) is contained
    in the kernel of some nonzero \(A\)-module homomorphism
    \[
        F:S^m\longrightarrow S.
    \]
    Every such homomorphism has the form
    \[
        F(z_1,\dots,z_m)
        =
        \varphi_1(z_1)+\cdots+\varphi_m(z_m),
    \]
    where
    \[
        \varphi_i\in \operatorname{End}_A(S),
    \]
    and not all \(\varphi_i\) are zero.

    Since \(N\subseteq \operatorname{Ker}(F)\), for every \(a\in A\) we have
    \[
        0
        =
        F(ax_1,\dots,ax_m)
        =
        \sum_{i=1}^m \varphi_i(ax_i).
    \]
    Because each \(\varphi_i\) is \(A\)-linear,
    \[
        \sum_{i=1}^m \varphi_i(ax_i)
        =
        \sum_{i=1}^m a\varphi_i(x_i)
        =
        a\left(\sum_{i=1}^m \varphi_i(x_i)\right).
    \]
    Taking \(a=1\), we obtain
    \[
        \sum_{i=1}^m \varphi_i(x_i)=0.
    \]
    In terms of the right \(D\)-module structure on \(S\), this is
    \[
        \sum_{i=1}^m x_i\cdot \varphi_i=0.
    \]
    Since \(x_1,\dots,x_m\) are right \(D\)-linearly independent, it follows that
    \[
        \varphi_1=\cdots=\varphi_m=0,
    \]
    contradicting the fact that \(F\neq 0\). Therefore
    \[
        N=S^m.
    \]
    Hence for any \((y_1,\dots,y_m)\in S^m\), there exists \(a\in A\) such that
    \[
        (ax_1,\dots,ax_m)=(y_1,\dots,y_m).
    \]
    This proves the density claim.

    Now let
    \[
        T\in \operatorname{End}_D(S).
    \]
    Since \(S\) is finite-dimensional over \(k\), it is also finite-dimensional as a
    right \(D\)-vector space. Choose a right \(D\)-basis
    \[
        s_1,\dots,s_n
    \]
    of \(S\). Applying the density claim to
    \[
        x_i=s_i,
        \qquad
        y_i=T(s_i),
    \]
    there exists \(a\in A\) such that
    \[
        as_i=T(s_i)
        \qquad
        (1\leq i\leq n).
    \]
    Since both \(L_a\) and \(T\) are right \(D\)-linear and they agree on the
    right \(D\)-basis \(s_1,\dots,s_n\), we get
    \[
        L_a=T.
    \]
    Thus every element of \(\operatorname{End}_D(S)\) lies in the image of \(\rho\).
    Therefore \(\overline{\rho}\) is surjective.

    Hence
    \[
        A/I\cong \operatorname{End}_D(S).
    \]
\end{proof}

\begin{remark}
    If one writes
    \[
        D=\operatorname{End}_A(S)
    \]
    instead of \(D=\operatorname{End}_A(S)^{\operatorname{op}}\), then the same statement
    is obtained by remembering that \(S\) is naturally a right
    \(\operatorname{End}_A(S)^{\operatorname{op}}\)-vector space via
    \[
        s\cdot \varphi=\varphi(s).
    \]
    In the split case \(D=k\), this distinction disappears.
\end{remark}

\begin{proposition}
    Let \(A\) be a finite-dimensional \(k\)-algebra, and let \(S\) be a simple
    left \(A\)-module. Then the following are equivalent:

    \begin{enumerate}
        \item \[
                  \operatorname{End}_A(S)=k\operatorname{id}_S.
              \]

        \item For every field extension \(K/k\), the scalar extension
              \[
                  S_K:=K\otimes_k S
              \]
              is a simple left \(A_K\)-module, where
              \[
                  A_K:=K\otimes_k A.
              \]
    \end{enumerate}
\end{proposition}

\begin{proof}
    Let
    \[
        D=\operatorname{End}_A(S).
    \]
    By Schur's lemma, \(D\) is a division algebra over \(k\).

    Assume first that
    \[
        D=k.
    \]
    Let
    \[
        I=\operatorname{Ann}_A(S).
    \]
    By the density theorem, since \(S\) is a finite-dimensional simple
    \(A\)-module, the natural map
    \[
        A/I\longrightarrow \operatorname{End}_D(S)
    \]
    is an isomorphism. Since \(D=k\), this becomes
    \[
        A/I\cong \operatorname{End}_k(S).
    \]
    After extending scalars from \(k\) to \(K\), we get
    \[
        A_K/(K\otimes_k I)
        \cong
        \operatorname{End}_K(K\otimes_k S).
    \]
    Thus the action of \(A_K\) on \(S_K=K\otimes_k S\) contains the full matrix
    algebra
    \[
        \operatorname{End}_K(S_K).
    \]
    Hence every nonzero \(A_K\)-submodule of \(S_K\) is stable under the full
    endomorphism algebra of \(S_K\), and therefore must be all of \(S_K\). Thus
    \(S_K\) is simple.

    Conversely, assume that for every field extension \(K/k\), the \(A_K\)-module
    \(S_K\) is simple. Since \(A\) is finite-dimensional and \(S\) is finite-dimensional,
    base change gives
    \[
        K\otimes_k \operatorname{End}_A(S)
        \cong
        \operatorname{End}_{A_K}(S_K).
    \]
    Take \(K=\overline{k}\). Then \(S_{\overline{k}}\) is simple by assumption, so
    Schur's lemma implies that
    \[
        \operatorname{End}_{A_{\overline{k}}}(S_{\overline{k}})
    \]
    is a division algebra. Therefore
    \[
        \overline{k}\otimes_k D
    \]
    is a finite-dimensional division algebra over the algebraically closed field
    \(\overline{k}\).

    But any finite-dimensional division algebra over an algebraically closed field
    is just the field itself. Hence
    \[
        \overline{k}\otimes_k D\cong \overline{k}.
    \]
    Taking dimensions over \(\overline{k}\), this implies
    \[
        \dim_k D=1.
    \]
    Therefore
    \[
        D=k,
    \]
    that is,
    \[
        \operatorname{End}_A(S)=k\operatorname{id}_S.
    \]
\end{proof}

\begin{remark}
    If \(k=\overline{k}\), then Schur's lemma implies that \(k\) is a splitting
    field for \(G\). Indeed, for every simple \(kG\)-module \(S\),
    \[
        \operatorname{End}_{kG}(S)=k\operatorname{id}_S.
    \]
\end{remark}

\begin{theorem}[Brauer]
    Let \(G\) be a finite group. If \(k\) contains a primitive
    \(|G|\)-th root of unity, then \(k\) is a splitting field for \(G\).
\end{theorem}

\begin{proof}
    Put
    \[
        n=|G|.
    \]
    Since \(k\) contains a primitive \(n\)-th root of unity, we have
    \[
        \operatorname{char}k\nmid n.
    \]
    Hence, by Maschke's theorem, \(kG\) is semisimple.

    Let \(\overline{k}\) be an algebraic closure of \(k\). It is enough to prove
    that every simple \(\overline{k}G\)-module is realizable over \(k\). Indeed, if
    every simple \(\overline{k}G\)-module has a \(k\)-form, then for every simple
    \(kG\)-module \(S\), the scalar extension
    \[
        \overline{k}\otimes_k S
    \]
    is again simple. Equivalently,
    \[
        \operatorname{End}_{kG}(S)=k,
    \]
    so \(k\) is a splitting field for \(G\).

    Let \(V\) be a simple \(\overline{k}G\)-module, and let
    \[
        \chi:G\longrightarrow \overline{k}
    \]
    be its character. For each \(g\in G\), the order of \(g\) divides \(n\), so
    \[
        g^n=1.
    \]
    Therefore the linear operator representing \(g\) on \(V\) satisfies
    \[
        X^n-1=0.
    \]
    Since \(k\) contains all \(n\)-th roots of unity and
    \(\operatorname{char}k\nmid n\), the polynomial \(X^n-1\) splits over \(k\)
    with distinct roots. Hence all eigenvalues of \(g\) on \(V\) lie in \(k\).
    It follows that
    \[
        \chi(g)\in k
    \]
    for all \(g\in G\). Thus every irreducible character of \(G\) over
    \(\overline{k}\) is \(k\)-valued.

    We now use Brauer's induction theorem. It says that every irreducible character
    \(\chi\) of \(G\) is an integral linear combination of characters induced from
    linear characters of subgroups. More precisely,
    \[
        \chi
        =
        \sum_i a_i\,\operatorname{Ind}_{H_i}^G(\lambda_i),
        \qquad a_i\in \mathbb{Z},
    \]
    where each \(H_i\leq G\), and each \(\lambda_i\) is a one-dimensional character
    of \(H_i\).

    Since every element of \(H_i\) has order dividing \(n\), the values of
    \(\lambda_i\) are \(n\)-th roots of unity. Hence
    \[
        \lambda_i(H_i)\subseteq k.
    \]
    Therefore each \(\lambda_i\) is afforded by a one-dimensional \(kH_i\)-module.
    Consequently,
    \[
        \operatorname{Ind}_{H_i}^G(\lambda_i)
    \]
    is afforded by a \(kG\)-module.

    Thus \(\chi\) is an integral linear combination of characters afforded by
    \(kG\)-modules.

    Now recall the standard Schur-index criterion: if \(\chi\) is an irreducible
    \(\overline{k}\)-character which is \(k\)-valued, then the Schur index
    \(m_k(\chi)\) is the greatest common divisor of the multiplicities with which
    \(\chi\) occurs in scalar extensions of \(kG\)-modules. In particular, if
    \(\chi\) is an integral linear combination of \(kG\)-characters, then
    \[
        m_k(\chi)\mid 1.
    \]
    Hence
    \[
        m_k(\chi)=1.
    \]
    Therefore \(V\) is realizable over \(k\). That is, there exists a simple
    \(kG\)-module \(S\) such that
    \[
        \overline{k}\otimes_k S\cong V.
    \]

    Since this holds for every simple \(\overline{k}G\)-module \(V\), every simple
    \(kG\)-module remains simple after extension of scalars to \(\overline{k}\).
    Equivalently,
    \[
        \operatorname{End}_{kG}(S)=k
    \]
    for every simple \(kG\)-module \(S\). Hence \(k\) is a splitting field for
    \(G\).
\end{proof}

\begin{proposition}
    Assume that
    \[
        \operatorname{char} k=p
    \]
    and that \(G\) is a finite \(p\)-group. Then the only simple \(kG\)-module is
    the trivial module \(k\).
\end{proposition}

\begin{proof}
    Let \(S\) be a simple \(kG\)-module. We show that \(S\) contains a nonzero
    \(G\)-fixed vector.

    Choose
    \[
        0\neq x\in S.
    \]
    Consider the finite subset
    \[
        X=\{\,gx\mid g\in G\,\}\subseteq S.
    \]
    Let \(S_0\) be the subgroup of the additive group of \(S\) generated by \(X\).
    Since \(\operatorname{char}k=p\), the additive group of \(S\) has exponent
    \(p\). Hence \(S_0\) is a finite abelian \(p\)-group.

    Moreover, \(S_0\) is stable under the action of \(G\). Indeed, if \(h\in G\)
    and \(gx\in X\), then
    \[
        h(gx)=(hg)x\in X.
    \]
    Thus \(G\) acts on the finite \(p\)-group \(S_0\).

    We now use the standard fact that if a finite \(p\)-group acts on a finite
    \(p\)-group, then the set of fixed points has cardinality congruent to the
    whole set modulo \(p\). Equivalently,
    \[
        |S_0^G|\equiv |S_0|\pmod p.
    \]
    Since \(S_0\) is a finite \(p\)-group, we have
    \[
        |S_0|\equiv 0\pmod p.
    \]
    Therefore
    \[
        |S_0^G|\equiv 0\pmod p.
    \]
    But \(0\in S_0^G\), so \(S_0^G\) is not empty. Since its cardinality is
    divisible by \(p\), it must contain an element different from \(0\). Hence
    there exists
    \[
        0\neq y\in S_0
    \]
    such that
    \[
        gy=y
        \qquad
        \text{for all }g\in G.
    \]

    Now consider the \(k\)-subspace
    \[
        ky\subseteq S.
    \]
    Since \(gy=y\) for all \(g\in G\), the subspace \(ky\) is stable under the
    action of \(G\). Hence \(ky\) is a nonzero \(kG\)-submodule of \(S\).

    Because \(S\) is simple, we must have
    \[
        S=ky.
    \]
    Therefore \(S\) is one-dimensional, and \(G\) acts trivially on \(S\). Hence
    \[
        S\cong k
    \]
    as \(kG\)-modules, where \(k\) is the trivial \(kG\)-module.
\end{proof}

\begin{remark}
    Thus, when \(\operatorname{char}k=p\) and \(G\) is a finite \(p\)-group, the
    group algebra \(kG\) has only one simple module up to isomorphism, namely the
    trivial module.
\end{remark}

\begin{theorem}
    Assume that
    \[
        \operatorname{char}k\nmid |G|,
    \]
    or equivalently, that \(|G|\) is invertible in \(k\). Then \(kG\) is
    semisimple. Hence, as a ring,
    \[
        kG
        \cong
        \bigoplus_{i=1}^r M_{n_i}(D_i),
    \]
    where each \(D_i\) is a finite-dimensional division algebra over \(k\).
\end{theorem}

\begin{proof}
    By Maschke's theorem, \(kG\) is semisimple. Since \(kG\) is finite-dimensional
    over \(k\), it is an Artinian semisimple ring. Therefore the Artin--Wedderburn
    theorem implies that
    \[
        kG
        \cong
        \bigoplus_{i=1}^r M_{n_i}(D_i),
    \]
    where each \(D_i\) is a division algebra finite-dimensional over \(k\).
\end{proof}

\begin{theorem}
    Assume further that \(k\) is a splitting field for \(G\). In particular, this
    holds if \(k=\overline{k}\). Let
    \[
        S_1,\dots,S_r
    \]
    be a complete set of pairwise non-isomorphic simple \(kG\)-modules, and set
    \[
        d_i=\dim_k S_i.
    \]
    Then, as a left \(kG\)-module,
    \[
        kG
        \cong
        S_1^{\oplus d_1}
        \oplus
        \cdots
        \oplus
        S_r^{\oplus d_r}.
    \]
    Consequently,
    \[
        |G|
        =
        d_1^2+\cdots+d_r^2.
    \]
    Moreover, as a \(k\)-algebra,
    \[
        kG
        \cong
        M_{d_1}(k)
        \oplus
        \cdots
        \oplus
        M_{d_r}(k).
    \]
\end{theorem}

\begin{proof}
    Since \(kG\) is semisimple, every simple module is projective and every
    \(kG\)-module is semisimple. Thus the regular module decomposes as
    \[
        kG
        \cong
        S_1^{\oplus m_1}
        \oplus
        \cdots
        \oplus
        S_r^{\oplus m_r}
    \]
    for some multiplicities \(m_i\).

    Because \(k\) is a splitting field, Schur's lemma gives
    \[
        \operatorname{End}_{kG}(S_i)=k.
    \]
    The Artin--Wedderburn decomposition therefore has the form
    \[
        kG
        \cong
        M_{d_1}(k)\oplus\cdots\oplus M_{d_r}(k),
    \]
    where \(d_i=\dim_k S_i\). As a left module over \(M_{d_i}(k)\), the algebra
    \(M_{d_i}(k)\) is a direct sum of \(d_i\) copies of its natural simple module
    \(k^{d_i}\). Hence
    \[
        kG
        \cong
        S_1^{\oplus d_1}
        \oplus
        \cdots
        \oplus
        S_r^{\oplus d_r}
    \]
    as a left \(kG\)-module.

    Taking \(k\)-dimensions gives
    \[
        |G|
        =
        \dim_k kG
        =
        \sum_{i=1}^r d_i\dim_k S_i
        =
        \sum_{i=1}^r d_i^2.
    \]
\end{proof}
